%!TEX root = ../main.tex
% ANEXO I — Repositorios del proyecto MarketLens

\chapter{Repositorios de código fuente y documentación}
\label{anx:repositorios}

Este anexo identifica los repositorios públicos asociados con MarketLens,
con el fin de facilitar la consulta del código fuente, las instrucciones
de ejecución y la documentación académica del proyecto.

\section{Repositorio del sistema MarketLens}
\label{anx:repositorio-software}

\noindent\textbf{URL:}
\url{https://github.com/TomasPuglieseUSAL/marketlens}

\noindent\textbf{Rama principal:} \texttt{main}.

\noindent\textbf{Estado:} repositorio público y activo al
16 de noviembre de 2026.

\noindent\textbf{Contenido:} implementación del prototipo,
organizada en un backend con Java y Spring Boot, una aplicación web
con React y TypeScript, una extensión experimental de navegador,
scripts auxiliares y archivos para su ejecución con Docker Compose.

\noindent\textbf{Persistencia y autenticación:}
PostgreSQL, Spring Security y tokens JWT.

\noindent\textbf{Reproducción:} el archivo \texttt{README.md} del
repositorio detalla los requisitos, las variables de configuración
local, los procedimientos de inicio con Docker Compose y las
instrucciones para cargar la extensión del navegador.
La instalación documentada es local y experimental; no constituye
un servicio de producción.

\noindent\textbf{Licencia:} no se identifica un archivo
\texttt{LICENSE} en la raíz de la versión consultada. La
disponibilidad pública facilita la revisión académica del código,
pero no implica por sí misma una licencia de reutilización.

\section{Repositorio del informe académico}
\label{anx:repositorio-documentacion}

\noindent\textbf{URL:}
\url{https://github.com/TomasPuglieseUSAL/MarketLens-Proyecto-Final}

\noindent\textbf{Rama principal:} \texttt{main}.

\noindent\textbf{Estado:} repositorio público al
16 de noviembre de 2026.

\noindent\textbf{Contenido:} fuentes de este informe en
\LaTeX, capítulos, archivos preliminares, bibliografía,
figuras y material complementario, organizados sobre
la plantilla institucional del Proyecto Final de Ingeniería
Informática de la Universidad del Salvador.

\noindent\textbf{Compilación:} el archivo \texttt{main.tex}
es el punto de entrada del documento. La plantilla utiliza
pdfLaTeX y las herramientas bibliográficas y de glosarios
indicadas en su documentación.
